\chapter{Advanced Volatility and Dependence}

\section{Beyond GARCH}

Chapter~31 covered GARCH, EGARCH and GJR-GARCH. This chapter treats
stochastic volatility, dynamic correlations, copulas, wavelets, and
self-exciting point processes.

\section{Stochastic Volatility}

GARCH models the conditional variance as a deterministic function of
past squared errors. Stochastic Volatility (Taylor, 1986) treats the
log-volatility as a latent AR(1) process:

\[
y_t = \exp(h_t/2) \varepsilon_t, \qquad
h_t = \mu + \phi(h_{t-1} - \mu) + \eta_t
\]

\begin{lstlisting}[caption={Stochastic Volatility}]
let n = 500
generate df y = 0.02 + rnormal(n, 0, exp(0.5*rnormal(n, 0, 1)))

let m_sv = sv(df, y, iter=500)
print(m_sv)
\end{lstlisting}

\texttt{sv} reports $\mu$, $\phi$ (persistence), $\sigma_\eta$ (vol-of-vol),
and the conditional volatility path.

\section{Dynamic Conditional Correlation}

DCC-GARCH (Engle, 2002) combines univariate GARCH(1,1) models with a
dynamic equation for the conditional correlation matrix. It is useful
for modelling time-varying dependence between returns.

\begin{lstlisting}[caption={DCC-GARCH}]
let m_dcc = dcc_garch(y1 ~ y2, df)
print(m_dcc)
\end{lstlisting}

The output includes the DCC parameters ($\alpha$, $\beta$), univariate
GARCH parameters, conditional volatilities, and dynamic correlation
matrices.

\section{Copulas}

\texttt{copula} estimates a bivariate dependence model via Sklar's theorem. It
separates marginal distributions from the copula that links them.

\begin{lstlisting}[caption={Copula}]
let m_cop = copula(y1 ~ y2, df, type="gaussian")
print(m_cop)
\end{lstlisting}

Available copulas include Gaussian, Clayton, Gumbel, and Frank. The
output reports Kendall's $\tau$, Spearman's $\rho$, and the copula
parameter.

\texttt{tv\_copula} extends this with time-varying dependence driven by a
GARCH-like dynamics for the copula parameter.

\begin{lstlisting}[caption={Time-varying copula}]
let m_tvc = tvcopula(y1 ~ y2, df, type="gaussian")
print(m_tvc)
\end{lstlisting}

\section{Wavelets}

Maximal Overlap Discrete Wavelet Transform (MODWT) decomposes a series
into time-frequency components without decimation.

\begin{lstlisting}[caption={MODWT}]
let m_modwt = modwt(df, y, scales=4)
print(m_modwt)
\end{lstlisting}

Useful for isolating cyclical components at different frequencies.

\section{Hawkes processes}

A Hawkes process is a self-exciting point process. Each event increases
the conditional intensity of future events by an amount that decays
exponentially.

\begin{lstlisting}[caption={Hawkes process}]
let m_hawkes = hawkes(df, time, T=100)
print(m_hawkes)
\end{lstlisting}

The branching ratio $\eta = \alpha / \beta$ tells how much each event
contributes to future events. $\eta < 1$ is required for stationarity.
